\chapter{Background and Related Work} \label{Chap2}

\section{North Wyke Farm Platform}

The North Wyke Farm Platform (NWFP), operated by Rothamsted Research in Devon, UK, comprises three
hydrologically isolated, intensively-monitored grazed grassland catchments -- referred to
throughout this dissertation as Towers 2, 4, and 9 -- each instrumented with an Eddy Covariance
(EC) flux tower since 2017--2018. The three catchments are managed as beef-cattle farmlets under a
shared research programme, and NWFP's broader research output, including EC-based flux studies of
CO\textsubscript{2} \cite{cardenas2022} and spatial variability in greenhouse gas fluxes at
comparable UK sheep-grazed pastures \cite{charteris2021spatial}, positions it as a long-running,
representative case study of intensively-managed UK grassland livestock systems rather than a
single, idiosyncratic site. NWFP has also been the subject of prior quantile-machine-learning work
on soil moisture \cite{oulaid2025} and animal-scale methane prediction from the platform's GreenFeed
system \cite{partridge2024}, and is explicitly named as a candidate site for a farm-scale digital
twin \cite{fakeye2024}. Full technical specifications of the dataset used in this dissertation --
catchment areas, sensor configurations, and coordinates -- are given in Chapter~\ref{Chap3}.

\section{Eddy Covariance Measurement of FCH4 and Other Gases}

\subsection{Gap-Filling Methodology and Its Difficulty}

The EC-ML literature for CH\textsubscript{4} flux prediction is dominated by tree-based ensemble
methods. Irvin et al.\ established a benchmark across 17 FLUXNET-CH4 wetland sites, finding Random
Forest and XGBoost superior to Artificial Neural Networks and penalised linear regression across
all biomes, with soil temperature the most important predictor, and documented that raw ML
uncertainty estimates were systematically underestimated across every site tested
\cite{irvin2021}. Kim et al.\ confirmed Random Forest dominance across five wetland and rice paddy
sites, adding that lagged environmental variables improve CH\textsubscript{4} prediction
substantially more than CO\textsubscript{2} does \cite{kim2020}. Zhu et al.\ extended these
findings to challenging ecosystems including UK managed pastures, validating Random-Forest-based
gap-filling over Marginal Distribution Sampling (MDS) as the recommended approach for precisely
the ecosystem type studied at NWFP \cite{zhu2023a}, and Kheradmand et al.\ confirm a similar
pattern for CO\textsubscript{2} gap-filling at grassland and cropland sites using tower-based EC
data \cite{kheradmand2023}. Across this literature, the models used are almost exclusively simple
ANN and tree-based ensembles -- a contrast this chapter returns to in
Section~\ref{sec:tabular-modelling}.

\subsection{Lack of Forecasting/Scenario-Projection Work in This Sector}

This literature is nonetheless answering a different question from the one this dissertation
asks. Gap-filling is a within-distribution interpolation task: a model reconstructs missing
observations inside a known regime, using both past and future context around the gap. Forecasting
is sequential extrapolation: a model predicts forward from real or previously-predicted history
alone, with no access to future context, and temporal dependencies become the dominant modelling
challenge rather than a secondary concern. Every gap-filling benchmark reviewed above answers the
former question; none addresses the latter. To the author's knowledge, no prior study has applied
ML-based forecasting to EC CH\textsubscript{4} flux at a managed temperate grassland -- this is the
central gap this dissertation addresses.

\section{Towards Digital Twins and the What-If Simulation Gap}

Agricultural digital twins are difficult to build precisely because agricultural systems are
living, weather-dependent, and non-stationary, unlike the controlled, deterministic environments
industrial digital twins are usually built for. Purcell and Neubauer's systematic review of
agricultural digital twins finds that what-if scenario simulation -- the capability that would
transform a digital twin from a monitoring tool into a decision-support system -- remains
critically absent across the field \cite{purcell2023state}. Fakeye et al.\ propose a three-tier
digital twin architecture specifically for NWFP and name CH\textsubscript{4} forecasting as a
missing module within it \cite{fakeye2024}. Four components motivate this dissertation's response
to that gap:

\begin{enumerate}
    \item \textbf{Forecasting} -- as established in Section 2.2.2, this gap exists independently
    of any digital-twin framing.
    \item \textbf{Scenario projection} -- the specific what-if capability Purcell and Neubauer
    identify as missing.
    \item \textbf{Interpretability} -- Buzacott et al.\ demonstrate via Shapley-value analysis
    that raw predictive accuracy can mask a highly entangled multi-driver dependency in
    CH\textsubscript{4} emissions \cite{buzacott2024}; Sharma et al.\ caution that models lacking
    site-specific calibration frequently miss episodic, high-emission events \cite{sharma2026};
    Fuchs et al.\ show that ignoring discrete management interventions leads to systematic
    under-prediction of peak emission fluxes \cite{fuchs2020}.
    \item \textbf{Uncertainty quantification} -- anchored on Irvin et al.'s finding that raw ML
    uncertainty is systematically underestimated \cite{irvin2021}: a model that performs well on
    average but misrepresents its own confidence is not safe to use for scenario-based
    decision-support.
\end{enumerate}

\section{Recent Developments in Tabular Modelling Approaches}
\label{sec:tabular-modelling}

Section 2.2.1 shows the EC-ML gap-filling literature is dominated by simple ANN and tree-based
ensembles. A newer generation of tabular and sequence architectures, not yet applied in this
domain, motivates the model roster benchmarked in Chapters~\ref{Chap4} and~\ref{Chap5}.

\subsection{Tabular Foundation Models}

TabPFN-family foundation models are pretrained on large volumes of synthetic tabular data and make
predictions via in-context learning at inference time, without any dataset-specific parameter
fitting. TabICL extends this approach to substantially larger tables via a two-stage
column-then-row attention architecture \cite{qu2025tabicl}, and TabICLv2 further improves
generalisation and pretraining efficiency \cite{qu2026tabiclv2}. Hoo et al.\ extend the underlying
TabPFN-v2 model to time-series forecasting by treating it as a tabular regression problem with
lightweight temporal featurisation \cite{hoo2026tabpfnts} -- the same family of models this
dissertation applies to CH\textsubscript{4} forecasting (Chapter~\ref{Chap5}) and scenario
projection (Chapter~\ref{Chap6}), and a family with a documented structural limitation under
genuine distribution shift, returned to in Chapter~\ref{Chap6}'s discussion of extrapolation risk.

\subsection{Sequence Architectures and the Attention-vs-Linear Debate}

Whether architectural sophistication translates into forecasting accuracy is itself an open
question in the wider time-series literature. Zeng et al.\ demonstrate across nine standard
long-term forecasting benchmarks that a set of embarrassingly simple linear models outperform
sophisticated Transformer-based architectures, attributing the failure to permutation-invariant
self-attention destroying temporal ordering information \cite{zeng2023}. Montero-Manso and Hyndman
further show that global forecasting methods -- a single model fit across many related series --
can match or exceed local, per-series methods without assuming similarity between series
\cite{montero2021global}, directly relevant to this dissertation's pooled and per-tower model
comparisons. Together, this literature cautions against assuming any one architecture family will
win in advance, motivating the identical-conditions benchmarking pipeline adopted throughout this
dissertation.

\section{Synthesis and Gap Statement}

Read together, this literature establishes that ecosystem-scale CH\textsubscript{4} forecasting at
a managed temperate grassland is unaddressed: gap-filling benchmarks are mature but answer a
structurally different question (Section 2.2); what-if scenario simulation is a named, open gap in
the agricultural digital twin literature (Section 2.3); and the newer tabular and foundation model
families capable of addressing both problems have not yet been applied in this domain
(Section~\ref{sec:tabular-modelling}). To the author's knowledge, no prior study has applied
ML-based forecasting to EC CH\textsubscript{4} flux at a managed temperate grassland. This
dissertation's specific combination -- tabular and foundation models, applied to forecasting,
scenario projection, interpretability, and uncertainty quantification together, at a representative
UK managed grassland site -- is the necessary and appropriate response to that gap.
